\documentclass[../../main.tex]{subfiles}

\begin{document}

\section{Functions and References}

Up until now, we had everything in our \verb|main| function.
However, you may have noticed that dumping everything in there is
not efficient nor scalable. This is where functions kick in.

\subsection{Functions}

As always, let's start with definition.

\begin{definition}{}{}
A reusable code block dedicated for a certain task is called a
\textbf{function} \cite{CS02-5}.
\end{definition}

Functions are specially useful as we can use it multiple times just
by calling them instead of copy and pasting every time we need them.
Functions in C++ generally have the following format.

\begin{basicbox}
\begin{lstlisting}[style=cppstyle]
outputDataType functionName(parameters) {
    code block
    return outputValue;
}
\end{lstlisting}
\end{basicbox}

Looking at a bare \verb|main| function, we can see that we have
\verb|int| since the \verb|main| function returns an integer value
like $0$. Since no arguments are taken in a bare \verb|main| function,
we have empty parenthesis. There are some functions that do not
return any value. In that case, we would be using \verb|void|.
Moreover to use a function, we need to call it. Below is a quick
example.

\begin{cppbox}{CS0201.26052-01}{CS0201.26052-01}
\begin{lstlisting}[style=cppstyle]
#include <iostream>

void helloWorld() {
    std::cout << "Hello, World!\n";
}

int main() {
    helloWorld();
    helloWorld();
    helloWorld();
    return 0;
}
\end{lstlisting}
\end{cppbox}

In the example above, we have our \verb|helloWorld| function that
prints the text ``Hello, World!'' to the console. The type is
\verb|void| since it does not return any values. By calling the
function three times in \verb|main|, we don't have to manually
write long texts every time!

\begin{outputbox}
\begin{lstlisting}[style=basic]
Hello, World!
Hello, World!
Hello, World!
\end{lstlisting}
\end{outputbox}

Below is a more practical use of simple functions for basic
mathematical operations.

\begin{cppbox}{CS0201.26052-02}{CS0201.26052-02}
\begin{lstlisting}[style=cppstyle]
#include <iostream>

int add(int a, int b) {
    return a + b;
}

int subtract(int a, int b) {
    return a - b;
}

int multiply(int a, int b) {
    return a * b;
}

double divide(double a, double b) {
    return a / b;
}

int mod(int a, int b) {
    return a % b;
}

int main() {
    std::cout << "1 + 2 = "    << add(1, 2)      << std::endl;
    std::cout << "3 - 4 = "    << subtract(3, 4) << std::endl;
    std::cout << "5 * 6 = "    << multiply(5, 6) << std::endl;
    std::cout << "7 / 8 = "    << divide(7, 8)   << std::endl;
    std::cout << "10 mod 9 = " << mod(10, 9)     << std::endl;

    std::cout << "((1 + 2) * 4) mod 10 = "
              << mod(multiply(add(1, 2), 4), 10) << std::endl;

    return 0;
}
\end{lstlisting}
\end{cppbox}

Each function will take in the values $a$ and $b$ to perform the
dedicated task. After declaring the functions, we can call them in
the \verb|main| function to display them. Many times it is better
to make the functions declare a value rather than print directly
to the console for reusability. As shown in the last example,
we can call a function as an input of different function.

\begin{outputbox}
\begin{lstlisting}[style=basic]
1 + 2 = 3
3 - 4 = -1
5 * 6 = 30
7 / 8 = 0.875
10 mod 9 = 1
((1 + 2) * 4) mod 10 = 2
\end{lstlisting}
\end{outputbox}

With the ideas of functions in mind, we can continue with another
important topic of references.

\subsection{References}

One of the primary goals of writing codes for competitive programming
is efficiency. Indeed most, if not all, codes written for competitions
are not maintained as the main purpose is to solve problems
efficiently \cite{CS02-3}. Moreover, it is important that your
program is fast enough to execute under the time limit.








\end{document}


\begin{cppbox}
\begin{lstlisting}[style=cppstyle]

\end{lstlisting}
\end{cppbox}

\begin{outputbox}
\begin{lstlisting}[style=basic]

\end{lstlisting}
\end{outputbox}
